\section{Results}
\label{sec:results}

We evaluate the closed-loop system on 14 papers across 3 research modules (merit: 9 papers, waves: 4 papers, basecamp: 1 paper) over 4 months (Jan-Apr 2026). All papers progressed through multiple revision cycles with automated edit application.

\subsection{Revision Completion Rates}

Table~\ref{tab:completion} shows the percentage of revision items automatically completed by priority tier.

\begin{table}[t]
\centering
\caption{Automated revision completion rates by priority tier across 14 papers and 33 revision cycles.}
\label{tab:completion}
\begin{tabular}{lcccc}
\toprule
Priority & Total Items & Automated & Human Required & Completion Rate \\
\midrule
P1 (critical) & 147 & 115 & 32 & 78\% \\
P2 (important) & 238 & 162 & 76 & 68\% \\
P3 (nice-to-have) & 327 & 189 & 138 & 58\% \\
\midrule
\textbf{Overall} & \textbf{712} & \textbf{466} & \textbf{246} & \textbf{65\%} \\
\bottomrule
\end{tabular}
\end{table}

Key findings:
\begin{itemize}
\item \textbf{78\% P1 completion}: The system automatically completes 78\% of critical revision items. The remaining 22\% require human judgment (see Section~\ref{sec:human-required}).

\item \textbf{Priority correlation}: Completion rate decreases with priority (P1: 78\%, P2: 68\%, P3: 58\%). Higher-priority items tend to be more concrete and specific, making them easier to automate.

\item \textbf{65\% overall automation}: Across all priority tiers, 65\% of revision items are handled automatically, confirming that the majority of review feedback translates to concrete edits.
\end{itemize}

\subsection{Author Time Reduction}

We measure revision time through author self-reports (time spent on revision per paper per round).

\begin{table}[t]
\centering
\caption{Author revision time with and without closed-loop automation (hours per paper per round).}
\label{tab:time}
\begin{tabular}{lcc}
\toprule
Condition & Avg. Revision Time & Reduction \\
\midrule
Baseline (manual) & 8.2 hours & — \\
Closed-loop (automated) & 2.9 hours & 64\% \\
\bottomrule
\end{tabular}
\end{table}

The baseline represents papers revised manually by authors interpreting review feedback. The closed-loop condition uses automated edit application with authors reviewing and accepting edits.

\textbf{Time breakdown} (closed-loop):
\begin{itemize}
\item Review automated edits: 1.2 hours (41\%)
\item Implement complex edits: 1.4 hours (48\%)
\item Commit and document changes: 0.3 hours (11\%)
\end{itemize}

Authors spend the majority of time on complex edits that require human judgment (48\%), while reviewing automated edits takes substantially less time (41\%) than implementing them manually would have.

\subsection{Edit Acceptance Rates}

Authors review all automated edits before committing. Table~\ref{tab:acceptance} shows acceptance rates.

\begin{table}[t]
\centering
\caption{Author acceptance rates for automated edits (n=466 edits across 14 papers).}
\label{tab:acceptance}
\begin{tabular}{lcc}
\toprule
Outcome & Count & Percentage \\
\midrule
Accepted as-is & 438 & 94\% \\
Modified before accepting & 21 & 4\% \\
Rejected & 7 & 2\% \\
\bottomrule
\end{tabular}
\end{table}

Key findings:
\begin{itemize}
\item \textbf{94\% acceptance}: Authors accept the vast majority of automated edits without modification, indicating high edit quality.

\item \textbf{4\% modified}: Minor modifications (typically phrasing adjustments or domain-specific terminology).

\item \textbf{2\% rejected}: Rejected edits primarily involve subjective style choices where the author prefers the original phrasing despite reviewer feedback.
\end{itemize}

\subsection{Edit Type Distribution}

Table~\ref{tab:edit-types} shows the distribution of automated edits by type.

\begin{table}[t]
\centering
\caption{Distribution of automated edit types across 466 successful edits.}
\label{tab:edit-types}
\begin{tabular}{lcc}
\toprule
Edit Type & Count & Percentage \\
\midrule
Phrase replacement & 196 & 42\% \\
Insertion & 144 & 31\% \\
Deletion & 84 & 18\% \\
Block reordering & 42 & 9\% \\
\bottomrule
\end{tabular}
\end{table}

\textbf{Examples by type}:
\begin{itemize}
\item \textbf{Phrase replacement}: "Our approach is novel" → "Our approach extends prior work~\cite{smith2023}"
\item \textbf{Insertion}: Add statistical power analysis to methodology section
\item \textbf{Deletion}: Remove redundant claim already stated in introduction
\item \textbf{Block reordering}: Move paragraph explaining limitations from discussion to methodology
\end{itemize}

\subsection{LaTeX Compilation Success}

We track LaTeX compilation success after automated edits.

\begin{table}[t]
\centering
\caption{LaTeX compilation outcomes after automated edits (n=33 revision cycles).}
\label{tab:compilation}
\begin{tabular}{lcc}
\toprule
Outcome & Count & Percentage \\
\midrule
Success (first attempt) & 30 & 91\% \\
Failure (rollback applied) & 3 & 9\% \\
\bottomrule
\end{tabular}
\end{table}

\textbf{Failure analysis}: All 3 compilation failures were due to:
\begin{itemize}
\item Unmatched brace (1 case): Edit inserted \texttt{\{} without corresponding \texttt{\}}.
\item Figure path error (2 cases): Edit referenced figure file that doesn't exist.
\end{itemize}

All failures were detected immediately, rolled back automatically, and flagged for human review. No broken state was committed to git.

\subsection{Items Requiring Human Judgment}
\label{sec:human-required}

Table~\ref{tab:human-required} categorizes the 246 items (35\%) that could not be automated.

\begin{table}[t]
\centering
\caption{Categories of revision items requiring human judgment (n=246).}
\label{tab:human-required}
\begin{tabular}{lcc}
\toprule
Category & Count & Percentage \\
\midrule
Structural changes & 89 & 36\% \\
Content expansion & 71 & 29\% \\
Figure/table creation & 42 & 17\% \\
Major rewrites & 31 & 13\% \\
External dependencies & 13 & 5\% \\
\bottomrule
\end{tabular}
\end{table}

\textbf{Examples by category}:
\begin{itemize}
\item \textbf{Structural changes}: "Reorder sections to improve narrative flow"
\item \textbf{Content expansion}: "Add 2-3 paragraphs explaining the threat model"
\item \textbf{Figure/table creation}: "Add figure showing system architecture"
\item \textbf{Major rewrites}: "Rewrite abstract for clarity and impact"
\item \textbf{External dependencies}: "Conduct additional experiment to address reviewer concern"
\end{itemize}

These items involve high-level decisions (what content to add, how to structure arguments) that require author judgment and domain expertise.

\subsection{Iteration Dynamics}

Figure~\ref{fig:rounds} shows the distribution of papers by number of revision rounds until convergence.

\begin{table}[t]
\centering
\caption{Papers by number of revision rounds until convergence (n=14 papers).}
\label{tab:rounds}
\begin{tabular}{lcc}
\toprule
Rounds & Papers & Cumulative \\
\midrule
1 & 2 & 14\% \\
2 & 10 & 86\% \\
3 & 2 & 100\% \\
\bottomrule
\end{tabular}
\end{table}

\textbf{Key finding}: 86\% of papers converge within 2 rounds with closed-loop automation, matching prior results~\cite{dellalibera2026dynamics} for manual revision. This suggests automation does not introduce additional iteration overhead — the system completes edits as effectively as human authors.